\section{Select Exercise Solutions}\label{sec:solutions}
  \subsubsection{Solution to \cref{ex:unitarity}}\label{sec:solutions:unitarity}
    Let $U$ be such that, for all
      states $\ket{\psi}$, $\norm{U\ket{\psi}} = \norm{\ket{\psi}}$. Using
      the definition of the norm, this gives us that $\abs{(U\ket{\psi})^\dagger
      U\ket{\psi}} = \abs{\bra{\psi} U^\dagger U \ket{\psi}} =
      \abs{\braket{\psi|\psi}}$. Since this
      should hold for any $\ket{\psi}$, we must have that $U^\dagger U =
      \idmatrix = U U^\dagger$; the latter equality follows from the fact that,
      for all $\ket{\varphi}$ and $\ket{\xi}$, 
      $\abs{\braket{\varphi|\xi}} = \abs{\braket{\xi|\varphi}}$.
      Thus, $U$ satisfies the definition of an invertible matrix, with the
      inverse $U^\dagger$. Further, it is easy to see that the opposite also
      holds: any unitary matrix must be norm-preserving. Therefore, the
      set of norm-preserving operations on $N$-dimensional Hilbert space
      equal to the set of $N\times N$ unitary matrices.

  \subsubsection{Solution to \cref{ex:hermitian-projectors}}\label{sec:solutions:hermitian-projectors}
    Let $\projector_\varphi$ be the projector that
      projects a state $\ket \psi$ onto the one-dimensional subspace spanned
      by the state $\ket\varphi$. Then,
      \[
        \projector^\dagger_\varphi = (\ket{\varphi} \bra{\varphi})^\dagger 
        = \ket{\varphi} \bra{\varphi} = \projector_\varphi\, .
      \]

    \noindent
    Since $\ket\varphi$ is arbitrary, and any projector can be written as a sum of
      such one-dimensional projectors, \emph{and} a sum of Hermitian matrices must itself be Hermitian: 
      \[
        \projector^\dagger = \sum_i \projector^\dagger_i = \sum_i \projector_i = \projector\, ,
      \]
      this suffices to show that any projector must be Hermitian.
